\section{\ref{sec:pain}장. 통증}

통증 센서, 두 전선, 척수의 관문, 하행 음량 노브, 감작, 주의 포획은 직접 측정되었다. 관문과 감작의 식은 이 책의 단순화된 모델이다. 기계가 경보 음량을 고르는 규칙은 가설이다.

{\footnotesize
\setlength{\tabcolsep}{3pt}\renewcommand{\arraystretch}{1.15}
\begin{longtable}{>{\raggedright}p{0.5cm}>{\raggedright}p{4.0cm}>{\raggedright}p{5.6cm}>{\raggedright}p{2.3cm}>{\raggedright\arraybackslash}p{1.7cm}}
\toprule
No. & 명제 & 실험과 측정한 것 & 출처 & 상태 \\
\midrule\endfirsthead
\toprule
No. & 명제 & 실험과 측정한 것 & 출처 & 상태 \\
\midrule\endhead
1 & 높은 문턱값: 통증 센서마다 가열 문턱값은 $41^\circ$에서 $46$--$48^\circ$까지 & 단일 섬유 기록. 원숭이 피부: 느린(C) 센서의 중앙 문턱값 $41^\circ$, 빠른 제2형 센서 $46^\circ$, 제1형은 $53^\circ$ 이상. 사람 피부: 압력에도 응답하는 C 센서 $40.7^\circ$, 압력에 응답하지 않는 것 $48^\circ$. & \cite{Treede1995heat, Weidner1999cfib} & 측정됨 \\
2 & 통증 센서는 촉각 센서보다 훨씬 약하게 적응하고, 가벼운 손상 뒤에는 더 민감해진다; 강한 가열 뒤의 약화는 한 유형에서 측정됨 & 피부의 가벼운 화상($50^\circ$, $100\,\text{s}$): $5$--$10$~분 뒤 사람의 통증 문턱값과 원숭이 C 센서의 문턱값이 $2$--$6^\circ$ 내려간다; 다른 피부 센서에는 이런 이동이 없다. 빠른 제2형 센서는 강한 가열 뒤 응답이 억제된다. 촉각 센서와의 적응 비교는 일반적인 추정으로, 여기서 별도 실험으로 뒷받침되지는 않는다. & \cite{LaMotte1982hyper, Treede1995heat} & 측정됨 \\
3 & 두 전선: 빠른 전선 $5$--$30$~m/s, 느린 전선 $0.5$--$2$~m/s; 도착 시간 $t=L/v$~\eqref{eq:paintime} & 전도 속도: 원숭이의 빠른 통증 센서는 평균 $15$와 $25$~m/s(두 유형); 사람의 C 센서는 $0.8$--$1.0$~m/s. $L=1$~m이면 수십 밀리초와 약 1초. & \cite{Treede1995heat, Weidner1999cfib} & 측정됨 \\
4 & 통증은 두 번 온다: 먼저 빠르고 정확하게, 다음에 둔하고 오래 & 사람 아래팔 가열에 대한 반응 시간: 온도가 오르면 $1100$에서 $400\ms$로; 털 있는 피부의 강한 가열은 이중 통증을 주고, 손바닥은 주지 않는다. 첫 통증은 $6$~m/s보다 빠른 섬유만 나를 수 있다. 잠복기로 보면 손바닥에 없는 빠른 제2형 센서가 이에 해당한다. 역할 분담(“어디”와 “얼마나 나쁜지”)은 해석. & \cite{CampbellLaMotte1983first, Treede1995heat} & 측정됨 \\
5 & 관문: 척수 출력 뉴런이 통증과 촉각을 받고, 억제성 중개 뉴런은 촉각에 흥분한다(그림~\ref{fig:gate}) & 관문 도식은 이론으로 제안되었다. 생쥐, 유전적 표지와 뉴런 집단 제거: 소마토스타틴을 가진 흥분성 뉴런은 기계적 통증에 필요하다; 다이노르핀을 가진 억제성 뉴런은 촉각 센서의 입력이 그 뉴런에 닿지 못하게 한다. 이 뉴런이 없으면 가벼운 접촉이 통증을 일으킨다. & \cite{MelzackWall1965, Duan2014} & 측정됨 \\
6 & 촉각은 통증을 누그러뜨린다 & 사람, 손의 레이저 통증 펄스: 동시에 접촉하면 통증 평가와 펄스 에너지 구별 능력이 낮아진다; 한 피부 분절 안에서 접촉 지점과 통증 지점 사이 거리가 멀수록 효과가 약해진다. & \cite{Mancini2014touch} & 측정됨 \\
7 & 관문 이론의 가정: 강한 통증이 스스로 억제성 중개 뉴런을 꺼서 관문이 활짝 열린다 & 본문에서 원래 이론의 가정으로 표시했다. 생쥐 연구는 관문이 촉각을 통증 채널로 들이지 않는 방식을 기술한다; 통증에 의한 중개 뉴런의 꺼짐은 거기서 보이지 않았다. & \cite{MelzackWall1965} & 가설 \\
8 & 관문~\eqref{eq:gate}: 촉각 없이 문턱값 $0.18$, 촉각이 있으면 $0.86$, $g=2$에서 $0.11$; $\nu=1$에서 촉각은 출력을 $1$에서 $0.25$로 낮추고, $\nu=2$에서는 효과가 없다; 촉각 자체는 통과하지 못한다(그림~\ref{fig:gatecurves}) & 본문의 가중치로 \texttt{models/\allowbreak pain\_\allowbreak gate.py}에 따라 계산하면 이 수치가 모두 재현된다. & \texttt{models/\allowbreak pain\_\allowbreak gate.py} & 모델 \\
9 & 뇌간에서 내려오는 하행 경로는 관문의 이득을 양쪽으로 바꾼다 & 쥐, 연수, 꼬리 가열 중 기록: 어떤 뉴런은 꼬리를 빼기 전에 방전하고(“켜짐”), 다른 뉴런은 조용해진다(“꺼짐”); 이 영역의 약한 자극은 반사를 억제한다. 총설: 같은 회로가 통증 전달을 촉진하기도 억제하기도 하며, 오피오이드는 이를 통해 작용한다. & \cite{Fields1983rvm, Fields2004} & 측정됨 \\
10 & 통증이 줄어들 것이라는 기대는 오피오이드 채널을 통해 통증을 누그러뜨린다 & 급성 통증이 있는 사람들: 기대로 통증이 줄어든 사람에서는 오피오이드 수용체 차단이 통증을 다른 사람들 수준으로 되돌렸다; 다른 사람들에서는 차단이 통증을 바꾸지 않았다. & \cite{Levine1978} & 측정됨 \\
11 & 뇌는 중단의 이득에 따라 경보 음량을 고른다 & 음량 선택 규칙을 측정한 실험은 없다. & --- & 가설 \\
12 & 감작: 손상 뒤 관문 출력이 커지고 문턱값이 떨어지며, 일부는 척수에서 일어난다; 손상 자극이 멈춘 뒤에도 유지된다 & 쥐: 피부 손상 뒤 굽힘 반사의 문턱값이 떨어지고 응답이 커진다; 전기생리학적 분석은 증가의 일부가 척수 자체에서 생김을 보여 준다. 손상 자극은 자극 자체보다 오래가는 감도의 장기 변화를 일으킨다. 정확한 회복 시간은 본문에 없다. & \cite{Woolf1983} & 측정됨 \\
13 & 감작~\eqref{eq:sensitization}은 양의 되먹임이다; 루프가 강하면 경보는 치유 뒤에도 걸린 채 남을 수 있다 & 모델~\eqref{eq:sensitization}에서 나온다(장 끝의 연습문제). 치유 뒤 지속되는 통증이 바로 이런 형태로 걸린 루프임을 보인 실험은 없다. & 장 안의 유도 & 가설 \\
14 & 통증은 음의 보상이며, 그것으로 하는 학습은 시간차 오차를 따른다 & 영상 장치, 사람, 통증 전조의 사슬: 배쪽 선조체와 앞쪽 섬엽 피질의 활동이 시간차 학습 신호와 일치한다. 원숭이: \nt{DOPA} 뉴런 일부는 불쾌한 공기 분사와 그 전조에 흥분하고, 다른 일부는 억제된다. & \cite{Seymour2004, Matsumoto2009} & 측정됨 \\
15 & 통증은 진행 중인 작업을 중단시킨다 & 사람, 전기 통증 자극과 소리 시험: 통증 시작 $100\ms$ 뒤의 시험 응답은 뚜렷이 느려지고, $1.5\,\text{s}$ 뒤에는 대조군과 차이가 없다. 총설은 이런 실험을 주의 포획 모델로 정리한다. & \cite{Crombez1996disrupt, EcclestonCrombez1999} & 측정됨 \\
16 & 회피 반사는 척수에서 닫힌다 & 쥐: 후각 깊은 층의 뉴런이 개별 근육의 회피 반사 공간 패턴을 되풀이한다; 경수에서 역방향으로 흥분시킬 수 없으므로 척수 중개 뉴런이다. & \cite{Schouenborg1995wr} & 측정됨 \\
\bottomrule
\end{longtable}}
